\begin{frame}{UCB = 연속적인 절절점}
  \begin{columns}[T]
    \begin{column}{0.48\textwidth}
      \centering
      \includegraphics[width=\linewidth]{ch02_2_ucb_bonus.png}
      \vskip0.3em
      {\footnotesize 왼쪽: 보너스가 $1/\sqrt{N}$ 로 감소.
      오른쪽: 추정치 최하(1.60)인 C가 보너스 1.92 덕분에
      1위(3.52), 순수 탐욕 1위 A(3.11)를 앞섬.}
    \end{column}
    \begin{column}{0.5\textwidth}
      \textbf{한눈에 들어오는 구조}:
      \begin{itemize}
        \item $Q_t$ 차이(예: A-B의 0.50)가 \textbf{보너스 차이
          (0.35)보다 크면} $\to$ UCB도 순수 탐욕과 같은 팔을
          고름(활용이 이김).
        \item $Q_t$ 차이가 \textbf{보너스 차이보다 작으면} $\to$
          ``덜 당겨본'' 팔을 고름(탐험이 이김). C vs A: $Q_t$ 차
          A가 0.90 앞, 보너스 차 C가 1.31 앞 $\to$ \textbf{탐험
          이김}.
      \end{itemize}
      \vskip0.3em
      \kb{UCB는 ``탐험을 0이나 1로만 하는'' 것이 아니라, 매 스텝
      보너스와 추정치 차이를 비교하는 \textbf{연속적인} 절절점.}
      $N$ 이 커지면 보너스가 $1/\sqrt{N}$ 감소 $\to$
      자연스럽게 $Q_t$ 지배 영역으로.
    \end{column}
  \end{columns}
\end{frame}
